\sectionguard
\section{Scope, Status, and Qualifications}

\subsection*{Question and thesis boundary}

This study answers one question: what competent ecclesiastical
authority has actually done about Guadalupe, in what words and with
what object; what the surviving documentary record will and will not
carry; how the historicity dispute stands among the scholars who
conduct it; what the beatification and canonization of Juan Diego
establish; what the documented examinations of the image do and do not
show; how the Nahuatl narrative reads as Nahuatl; and what all of it
obliges. It is a source-audited study instrument,
ecclesial-judgment-first in method: not a devotional narrative, a
critical edition, a magisterial act, a materials-science report, a
Nahuatl philological study, or a history of Mexico. \emph{It offers no
judgment of its own on whether the events of December 1531 occurred},
and its refusal to do so is a finding about the state of the evidence
and about this study's own competence, not a verdict against either
side.

\subsection*{Corpus and exclusions}

\emph{Included.} The acts of competent authority located and listed at
\S1.2, each at its own object; the printed witnesses of 1648 and 1649,
with the 1649 imprint read and collated at the page-image level for
the passages quoted; the seventeenth- and eighteenth-century
apologetic sequence at the level of the 1785 compilation; the 1666
depositions at the level of the published descriptions; the shrine's
own published inventory of sixteenth-century documents, reproduced at
the shrine's level with ceilings assigned; the silence problem as
Garc\'ia Icazbalceta's 1883 letter states it, read in full; the
\work{Informaci\'on} of 1556 through O'Gorman's analysis and appendices,
read at length; the named modern positions in the historicity dispute
at the levels recorded in \S4.3; the cause of Juan Diego from the 1981
petition to the 2002 canonization; the documented examinations of the
image in 1666, 1751, 1787 and 1982, and the crown episode of
1887--1895; a traceability audit of seven widely circulated claims
about the image; the Nahuatl text's rhetoric read directly in the 1649
imprint; the reception chain from 1555 to Leo~XIV's message of 5
February 2026; and the doctrinal and procedural frame, including the
2024 DDF norms.

\emph{Excluded or bounded.} No adjudication of the historicity
question. No identification of any artist. No assertion about the
physical origin or manufacture of the image. No re-diagnosis of the
cure recognized in 2001. No project translation of the
\work{Nican mopohua}, and no reproduction of the received Nahuatl
prayer printed in the 1649 tract. No Nahuatl philological conclusion
about the text's date. No characterization of the motives of any named
living or recent person. No claim of a comprehensive registry of
apparitions: this is a single-event study. No reproduction of any
complete document.

\emph{Not consulted.} S\'anchez 1648; Becerra Tanco; Florencia; the
\work{Informaciones jur\'idicas} of 1666 in any edition; the
\work{Informaci\'on} of 1556 in manuscript or in any of its three
printed editions; Sahag\'un's \work{Historia general} in the Florentine
Codex or any critical edition; Motolin\'ia, Mendieta, Las Casas, the
Garc\'es letter, and the 1569--1570 \work{Descripci\'on del
Arzobispado}; Cabrera 1756 and Bartolache 1790 beyond their catalogue
records; the Sol Rosales and Guadarrama conservation reports as
printed in 2005; Callahan 1981; the \term{C\'odice Escalada} and every
publication about it; the books of Poole, Sousa--Poole--Lockhart,
Brading, Burkhart, von Wobeser, Gonz\'alez Fern\'andez and Ch\'avez;
the acts of the Fifth Mexican Provincial Council; the Roman acta of
the 1890--1894 Office cause; the 1894 Roman printing of the new Office;
and \work{Verbum Domini}.

\subsection*{Geography, chronology, and currentness}

Subject geography: Tepeyac and Mexico City; Cuauhtitl\'an; the
Archdiocese of Mexico; Rome. Subject range: 1531 (the events as
narrated) and 1555--1556 (the earliest securely documented cult) to
the acts current at the as-of date. \emph{As-of date: 2026-07-25.} All
mutable claims---the Dicastery for the Doctrine of the Faith's
documents index, the Holy See's document trees, the shrine's own
published pages, and the most recent papal reception---were checked on
that date. No act determining the supernatural origin of the events
was located, and no later act superseding any of the acts inventoried
was located. Stable historical facts are not given artificial
currentness. Every negative result in this study is a bounded search,
correctable by better evidence, never a guarantee.

\subsection*{Taxonomy and terminology}

Kept distinct throughout: \emph{received narrative} (the content of
the 1649 tract and its descendants); \emph{documentary witness} (a
dated text, at the level at which it was read); \emph{competent act}
(a papal, dicasterial, conciliar or diocesan act, at its own object);
\emph{liturgical concession} (Office, Mass, feast, calendar
inscription); \emph{honour} (coronation, basilica dignity, patronage);
\emph{papal teaching and piety} (exhortation, homily, message,
pilgrimage); \emph{hagiographical presentation} (the Holy See's
biography of a saint, in its own genre); \emph{scholarly position}
(reported at the level recorded at \S4.3); \emph{material finding}
(what an examination reported, at the level of the publication
reached); and \emph{editorial synthesis}. ``Reported'' is the default
register for the events of 1531 and for everything attributed to the
apparition. \term{Fides humana} names the prudent human assent invited
by approved private revelation. The six determinations of the 2024 DDF
norms are current categories and are applied to no legacy act.

\subsection*{Source hierarchy and method}

Priority: operative texts in identified official witnesses---\work{Acta
Apostolicae Sedis}, \work{Acta Sanctae Sedis}, \work{Notitiae}, and
the Holy See's web texts, with the \work{AAS}/\work{ASS} pages located
and read at first hand in the Holy See's own scans and the 1999
\work{Notitiae} decree read from rendered page images; then the 1649
imprint, read in the John Carter Brown Library digitization with the
quoted passages collated against page images and, for the Nahuatl
passages, against the independent 1887 printing by Agust\'in de la
Rosa; then open-access peer-reviewed and university-press scholarship
(O'Gorman; Cuadriello) at its own level; then public-domain
nineteenth- and eighteenth-century printed sources at their own level;
then institutional self-description (the shrine's document page) marked
as such; then reported accounts (press, reference works,
postulation-side accounts), always marked as reported.

Quotations are exact from the named witness. Latin, Spanish and
Nahuatl are quoted in the original where the wording is operative,
with editorial working translations that are visibly commentary. No
project or AI rendering is offered as a received translation. No text
is printed for vocal prayer anywhere in this study.

\emph{Assisted verification, recorded.} Two research passes---on the
Holy See's acts and on the textual witnesses of the \work{Nican
mopohua}---were performed by assisting agents of the same AI
contributor working under instructions to report only what they had
actually retrieved. Every load-bearing result of those passes that is
used in the publication was then independently re-retrieved and read
by the author of this study, with two exceptions, which are marked in
place: the two \work{Notitiae} loci for the 2002 and 2008 states of
the General Roman Calendar (\S8.4), and the enumeration of continuing
papal celebrations (\S8.4).

\subsection*{Qualifications and unresolved items}

(1)~The Latin text of Benedict~XIV's 1754 act was not reached in any
edition; the incipit \work{Non est equidem} is attested only in the
1785 Guadalupan compilation, and an index anomaly in the available
Mechelen \work{Bullarium} volume is recorded at \S8.2. Nothing about
the brief's wording is asserted. (2)~Leo~XIII's brief of 8 February
1887, the Sacred Congregation of Rites decree of 6 March 1894, and any
contemporary Roman notice of the 1895 coronation were not found in the
\work{Acta Sanctae Sedis} volumes searched; they are reported from
O'Gorman. (3)~No act declaring Our Lady of Guadalupe patroness of Latin
America was found in \work{AAS} 1910 or 1911, and no Guadalupan act of
Pius~XII was found in \work{AAS} 1946. (4)~No Latin canonization
formula for Juan Diego was located in \work{AAS} 94 (2002); the
consistorial decree there printed assigns his canonization to 30 July
2002, against the volume's own index and the dated homily. (5)~The
rank of the 9 and 12 December celebrations in the General Roman
Calendar is not asserted; the 1999 decree's \term{gradu festi} governs
the calendars of the Americas. (6)~The \work{Nican mopohua}'s date and
authorship are unresolved and are not adjudicated. (7)~The
\term{C\'odice Escalada} is reported only as a disputed object; no
primary publication of any examination was reached. (8)~The 1990s
objections are reported from press and reference accounts; no letter
was read, and two attempted retrievals returned HTTP~403. (9)~Callahan
1981 and the 1982 conservation reports were not read; the latter are
used at Cuadriello's level. (10)~Seven popular claims about the image
could not be traced to a reachable source and are recorded as such at
\S6.3; ``not traced'' means this study could not reach a source, not
that none exists. (11)~The identity of whoever removed the painted
crown between 1887 and 1895 is not established and is not asserted.
(12)~Miguel Le\'on-Portilla's \work{Tonantzin Guadalupe} (2000) was
verified as to bibliographic identity through a Library of Congress
record but was not read, and his position is not summarized in this
study. (13)~No independent Mariological, historical, canonical,
archival, Nahuatl-philological, liturgical, art-historical, conservation,
or materials-science review has occurred; no ecclesiastical review or
approval is claimed or implied.

\subsection*{Rights}

Project prose is original and intended for CC~BY~4.0. Quotations remain
under their own status. Public domain in text: the 1649
\work{Huei tlamahui\c{c}oltica}; the 1785 \work{Colecci\'on de obras y
opusculos}; Cabrera 1756; Bartolache 1790; Agust\'in de la Rosa 1887;
Garc\'ia Icazbalceta's \work{Carta} in its 1896 printing and in the
Biblioteca Virtual Miguel de Cervantes transcription; the
sixteenth-century testimony transcribed within O'Gorman's edition. The
John Carter Brown Library states that its digital-collection images are
made available under CC~BY~4.0 with the credit line ``Courtesy of the
John Carter Brown Library''; the Library of Congress states that it is
unaware of copyright or other restrictions in its World Digital Library
copy of the same 1649 work. Scan images from those and other
repositories were consulted for collation and are not redistributed
here. \work{Acta Apostolicae Sedis}, \work{Acta Sanctae Sedis},
\work{Notitiae}, and Holy See web texts are quoted within ordinary
scholarly limits with citation. Edmundo O'Gorman, \work{Destierro de
sombras}, is \copyright\ Universidad Nacional Aut\'onoma de M\'exico,
Instituto de Investigaciones Hist\'oricas, whose digital edition
authorizes non-lucrative reproduction provided the work is not mutilated
and the full source and its electronic address are cited; both
conditions are observed. Jaime Cuadriello's article in \work{Anales del
Instituto de Investigaciones Est\'eticas} is open access at SciELO
M\'exico and is quoted briefly with citation. The Bas\'ilica de Santa
Mar\'ia de Guadalupe's institutional pages are \copyright\ their
publisher and are used, briefly and with citation, as institutional
self-description. Sousa, Poole and Lockhart's \work{The Story of
Guadalupe} (Stanford University Press and UCLA Latin American Center
Publications, 1998) is in copyright, was not read for this study, and
is cited by identity only; no part of its translation is reproduced.
Online availability is not treated as a reuse licence; no complete
document is reproduced; and no received prayer text is printed anywhere
in this study.
